\kbvcontinued ಸಂಸ್ಕೃತಿಯನ್ನು ಉಳಿಸಿ ಬೆಳೆಸುತ್ತಾಳೆ. ಈ ಅಂಶಕ್ಕೆ ಒತ್ತುಕೊಡುವ ಉದ್ದೇಶದಿಂದ ಗಾದೆಯ ಉತ್ತರಾರ್ಧದಲ್ಲಿ “ಒಲೆಗೊಂದು ಕೊರಡು” ಎಂಬ ಮಾತನ್ನು ಸೇರಿಸಿದೆ. ಒಲೆಯಲ್ಲಿ ಸದಾಕಾಲ ಬೆಂಕಿ ಇದ್ದೇ ತೀರಬೇಕು ; ಇದೊಂದು ದೃಷ್ಟಾಂತ. ಹೀಗೆ ಸ್ಪಷ್ಟೀಕರಣ ನೀಡಬೇಕಾದ ವಿಷಯಕ್ಕೆ ಮತ್ತೊಂದು ವಿಷಯದ ದೃಷ್ಟಾಂತ ಪ್ರತಿಬಿಂಬದಂತೆ ಕಾಣುವಂತಿದ್ದರೆ ಅರ್ಥಸಾದೃಶ್ಯ ಸಮಂಜಸವಾಗಿರುವುದು.\par

\indent ಈ ರೀತಿಯಾಗಿ ಎರಡು ವಾಕ್ಯಗಳು ಅಥವಾ ವಾಕ್ಯಾಂಶಗಳು ಅರ್ಥಸಾದೃಶ್ಯದಿಂದ ಬಿಂಬ ಪ್ರತಿಬಿಂಬದಂತೆ ಕಂಡುಬಂದರೆ, ಅಂತಹ ಅಲಂಕಾರವೆನಿಸುವುದು.\par

\noindent\hspace*{1.2em}\kbvlabel{ಉದಾ}\hspace{0.6em}ಊಟ ಬಲ್ಲವನಿಗೆ ರೋಗವಿಲ್ಲ; ಮಾತು ಬಲ್ಲವನಿಗೆ ಜಗಳವಿಲ್ಲ.\par

\kbvsechead{(೫)\hspace{0.45em}ಅರ್ಥಾಂತರನ್ಯಾಸಾಲಂಕಾರ :}

\indent “ಮಹಾಭಾರತದ ಯುದ್ಧದಲ್ಲಿ “ಸತ್ಯವೇ ದೇವರು; ಸತ್ಯಪಾಂಡವರಿಗೆ ಜಯ ಲಭಿಸಿತು; ಮಾರ್ಗವನ್ನೇ ಅನುಸರಿಸಿ ನಡೆದ ಸತ್ಯಧರ್ಮಗಳನ್ನು ನಂಬಿದವರಿಗೆ ಗಾಂಧೀಜಿಯವರು ಮಹಾತ್ಮರೆನಿಸಿ ಜಯ ಕಟ್ಟಿಟ್ಟಿದವರು.” ಇಲ್ಲಿ ಕೊಂಡು ಲೋಕಪ್ರಸಿದ್ಧರಾದರು” ಎರಡು ವಾಕ್ಯಗಳಿವೆ. ಮೊದಲನೆಯ ವಾಕ್ಯದಲ್ಲಿ ಕಂಡುಬರುವ ವಿಷಯ ಕಂಡುಬರುವ ಲೋಕೋಕ್ತಿಯನ್ನು ಎರಡನೆಯ ವಾಕ್ಯದಲ್ಲಿ ಕಂಡುಬರುವ ಲೋಕೋಕ್ತಿಯು `ಬರುವ ದೃಷ್ಟಾಂತವು ಸಮರ್ಥಿಸುವುದು, ವುದು.\par

\indent ಈ ಪ್ರಕಾರವಾಗಿ, ಒಂದು ವಿಶೇಷ ವಾಕ್ಯವನ್ನು ಸರ್ವಕಾಲಕ್ಕೂ ಅನ್ವಯಿಸುವ ಹೇಳಿಕೆಯಿಂದಾಗಲೀ ಅಥವಾ ಸರ್ವಕಾಲಕ್ಕೂ ಅನ್ವಯಿಸುವ ಹೇಳಿಕೆಯನ್ನು ಒಂದು ವಿಶೇಷ ವಾಕ್ಯದಿಂದಾಗಲೀ ಸಮರ್ಥಿಸುವ ತೆರನಾದ ಅರ್ಥಾಲಂಕಾರವನ್ನು ಅರ್ಥಾಂತರನ್ಯಾಸಾಲಂಕಾರವೆನ್ನುವರು.\par
